%% ============================================================
%%  Дипломная (бакалаврская) работа
%%  «Детекция цвета, номеров, марки машины
%%   по материалам видеокамеры»
%%  МФТИ, 2025
%%  Компилируется: pdfLaTeX
%% ============================================================
\documentclass[14pt,a4paper]{report}

%% ── Кодировка и язык ──────────────────────────────────────
\usepackage[T2A]{fontenc}
\usepackage[utf8]{inputenc}
\usepackage[russian,english]{babel}

%% ── Поля страницы ─────────────────────────────────────────
\usepackage[left=30mm,right=15mm,top=20mm,bottom=20mm]{geometry}

%% ── Математика ────────────────────────────────────────────
\usepackage{amsmath,amssymb,mathtools}

%% ── Графика и таблицы ─────────────────────────────────────
\usepackage{graphicx}
\usepackage{booktabs}
\usepackage{array}
\usepackage{tabularx}
\usepackage{multirow}
\usepackage{colortbl}
\usepackage[table]{xcolor}

%% ── TikZ для схем ─────────────────────────────────────────
\usepackage{tikz}
\usetikzlibrary{shapes.geometric, arrows.meta, positioning,
                fit, backgrounds, calc, chains}

%% ── Оформление ────────────────────────────────────────────
\usepackage{indentfirst}
\usepackage{setspace}
\usepackage{enumitem}
\usepackage{caption}
\usepackage{subcaption}
\usepackage{float}
\usepackage{fancybox}
\usepackage{tcolorbox}

%% ── Гиперссылки ───────────────────────────────────────────
\usepackage[hidelinks,unicode]{hyperref}

%% ── Стиль заголовков глав ─────────────────────────────────
\usepackage{titlesec}
\titleformat{\chapter}[block]
  {\normalfont\Large\bfseries}{\thechapter}{1em}{}
\titlespacing*{\chapter}{0pt}{-10pt}{20pt}

%% ── Межстрочный интервал 1.5 ──────────────────────────────
\onehalfspacing

%% ── Нумерация по главам ───────────────────────────────────
\numberwithin{equation}{chapter}
\numberwithin{figure}{chapter}
\numberwithin{table}{chapter}

%% ── Цвета для схем ────────────────────────────────────────
\definecolor{bluenode}{RGB}{52,101,164}
\definecolor{greennode}{RGB}{30,130,76}
\definecolor{orangenode}{RGB}{200,100,20}
\definecolor{graynode}{RGB}{100,100,110}
\definecolor{lightblue}{RGB}{210,228,255}
\definecolor{lightgreen}{RGB}{210,240,215}
\definecolor{lightorange}{RGB}{255,230,200}
\definecolor{lightgray}{RGB}{235,235,240}
\definecolor{rowgray}{RGB}{245,245,248}

%% ── TikZ стили ────────────────────────────────────────────
\tikzset{
  block/.style={
    rectangle, rounded corners=4pt, minimum width=3cm,
    minimum height=1cm, text centered, text width=3.2cm,
    draw, font=\small
  },
  blk/.style={block, fill=lightblue,  draw=bluenode,   thick},
  grn/.style={block, fill=lightgreen, draw=greennode,  thick},
  org/.style={block, fill=lightorange,draw=orangenode, thick},
  gry/.style={block, fill=lightgray,  draw=graynode,   thick},
  arr/.style={-{Stealth[length=6pt]}, thick},
  darr/.style={arr, dashed},
  lbl/.style={font=\footnotesize, midway, right=2pt}
}

% ─────────────────────────────────────────────────────────────
\begin{document}

%% =====================  ТИТУЛЬНЫЙ ЛИСТ  =====================
\begin{titlepage}
\begin{center}
{\small
Федеральное государственное автономное образовательное\\
учреждение высшего образования\\[4pt]
«Московский физико-технический институт\\
(национальный исследовательский университет)»
}

\vspace{5cm}

{\large\bfseries
«Детекция цвета, номеров, марки машины\\
по материалам видеокамеры»
}\\[1.2em]
{\normalsize (бакалаврская работа)}
\end{center}

\vspace{4cm}

\hfill
\begin{minipage}{0.42\textwidth}
  \textbf{Студент:}\\
  Клименко Виталий Евгеньевич\\[1.8em]
  \rule{\linewidth}{0.4pt}
\end{minipage}

\vspace{2em}

\hfill
\begin{minipage}{0.42\textwidth}
  \textbf{Научный руководитель:}\\
  Назаров Алексей Николаевич\\[1.8em]
  \rule{\linewidth}{0.4pt}
\end{minipage}

\vfill
\begin{center}
Москва\\
2025
\end{center}
\end{titlepage}

%% =====================  АННОТАЦИЯ  ==========================
\chapter*{Аннотация}
\addcontentsline{toc}{chapter}{Аннотация}

Выпускная квалификационная работа посвящена исследованию и разработке программного
комплекса для автоматического определения характеристик транспортных средств по
видеопотоку с камер наблюдения: детекции автомобиля в кадре, локализации
номерного знака, распознавания символов, а также классификации цвета и типа кузова.

Предлагается конвейерная архитектура на основе методов глубокого обучения:
YOLOv8n~--- для детекции ТС и номерных знаков, ResNet18~--- для классификации
цвета (9 классов) и типа кузова (6 классов), EasyOCR~--- для оптического
распознавания символов кириллицы.

Обучение и тестирование проводились на открытых датасетах (VCoR, CompCars,
Roboflow License Plates). Оценка качества выполнялась по метрикам
mAP@0.50, Accuracy top-1/3, F1 per class и CER. Проведено сравнительное
исследование ResNet18 и MobileNetV2. Результатом является работоспособный
прототип системы видеоаналитики транспорта.

%% =====================  СОДЕРЖАНИЕ  =========================
\tableofcontents

%% =============================================================
%%  ГЛАВА 1 — ВВЕДЕНИЕ
%% =============================================================
\chapter{Введение}

\section{Актуальность задачи}

В настоящее время системы видеонаблюдения широко применяются для мониторинга
дорожного движения, обеспечения безопасности и контроля доступа на объекты.
Объём видеоданных, поступающих с камер, постоянно растёт, что делает ручную
обработку информации неэффективной и трудоёмкой~\cite{abramov2019}.

Практический интерес представляет автоматическое определение характеристик
автомобиля: государственного регистрационного знака, марки, модели и цвета.
Эти данные применяются в задачах контроля проезда, поиска транспортных средств
и учёта трафика~\cite{survey2021}. Однако качество видеозаписей с камер
наблюдения часто ограничено низким разрешением, изменениями освещённости,
погодными условиями и движением объектов.

\section{Постановка задачи}

Задача анализа транспортных средств по видеоданным представляется как
последовательность следующих этапов:

\begin{enumerate}[nosep]
  \item детекция автомобиля в кадре;
  \item локализация номерного знака внутри области автомобиля;
  \item оптическое распознавание символов номера;
  \item классификация визуальных атрибутов: цвет и тип кузова.
\end{enumerate}

Формально результат детекции объектов на изображении:
\begin{equation}
  D = \bigl\{(b_i,\, c_i,\, p_i)\bigr\}_{i=1}^{N},
\end{equation}
где $b_i = (x_i, y_i, w_i, h_i)$~--- ограничивающий прямоугольник,
$c_i$~--- класс объекта, $p_i$~--- уверенность детекции.

Качество локализации оценивается через Intersection over Union:
\begin{equation}
  \mathrm{IoU} = \frac{S(B_{\mathrm{pred}} \cap B_{\mathrm{gt}})}
                      {S(B_{\mathrm{pred}} \cup B_{\mathrm{gt}})}.
\end{equation}

\section{Обзор методов детекции транспортных средств}

\subsection{Сравнение подходов к детекции}

Эволюция методов детекции объектов прошла несколько этапов.
Таблица~\ref{tab:detectors} суммирует основные подходы.

\begin{table}[H]
\centering
\caption{Сравнение методов детекции объектов}
\label{tab:detectors}
\rowcolors{2}{rowgray}{white}
\begin{tabularx}{\textwidth}{lXXcc}
\toprule
\textbf{Метод} & \textbf{Принцип} & \textbf{Достоинства} & \textbf{mAP} & \textbf{FPS} \\
\midrule
HOG+SVM       & Ручные признаки, скользящее окно   & Прост в реализации, нет GPU  & низкий & 5--15 \\
Haar-каскады  & Каскад слабых классификаторов      & Очень быстрые на CPU          & низкий & 30+ \\
Faster R-CNN  & Двухэтапная сеть, RPN+cls          & Высокая точность              & высокий & 5--15 \\
SSD           & Одноэтапная, anchor-based          & Баланс скорости и точности    & средний & 46+ \\
\textbf{YOLOv8n} & \textbf{Одноэтапная, anchor-free} & \textbf{Скорость + точность} & \textbf{высокий} & \textbf{100+} \\
\bottomrule
\end{tabularx}
\end{table}

\subsection{Двухэтапные детекторы (Faster R-CNN)}

Detectors семейства Faster R-CNN работают в два этапа: Region Proposal Network
генерирует кандидаты, затем классификационная голова уточняет координаты.
Подход обеспечивает высокую точность, но требователен к вычислительным ресурсам
и не обеспечивает скорость, достаточную для обработки видео в реальном времени.

\subsection{Одноэтапные детекторы (YOLO)}

YOLOv8~\cite{survey2021} решает задачу детекции за один проход по сети.
Anchor-free архитектура, CSPDarkNet backbone и FPN neck обеспечивают
работу на GPU со скоростью 100+ FPS при конкурентоспособной точности.

\noindent\textbf{Обоснование выбора:} YOLOv8n оптимален для видеопотока
в реальном времени и совместим с форматами ONNX/TensorRT.

\section{Обзор методов автоматического распознавания номерных знаков (ANPR)}

Типовая ANPR-система включает несколько последовательных этапов,
схема которых представлена на рисунке~\ref{fig:anpr_flow}.

\begin{figure}[H]
\centering
\begin{tikzpicture}[node distance=0.7cm, every node/.style={font=\small}]
  \node[blk, text width=2.6cm] (cam)   {Видеокадр\\с камеры};
  \node[blk, text width=2.6cm, right=0.8cm of cam] (vdet)  {Детекция\\автомобиля};
  \node[blk, text width=2.6cm, right=0.8cm of vdet](pdet)  {Локализация\\номерного знака};
  \node[blk, text width=2.6cm, right=0.8cm of pdet](prep)  {Предо-\\бработка};
  \node[blk, text width=2.6cm, right=0.8cm of prep](ocr)   {OCR\\символов};
  \node[grn, text width=2.6cm, right=0.8cm of ocr] (post)  {Пост-\\обработка};

  \draw[arr] (cam)  -- (vdet);
  \draw[arr] (vdet) -- (pdet);
  \draw[arr] (pdet) -- (prep);
  \draw[arr] (prep) -- (ocr);
  \draw[arr] (ocr)  -- (post);

  \node[gry, text width=2.2cm, below=0.9cm of pdet] (aug)
        {Перспективное\\выравнивание};
  \draw[darr] (pdet) -- (aug);
  \draw[darr] (aug)  -- (prep);
\end{tikzpicture}
\caption{Типовая схема ANPR-системы}
\label{fig:anpr_flow}
\end{figure}

\subsection{Сравнение ANPR-подходов}

\begin{table}[H]
\centering
\caption{Сравнение подходов к ANPR}
\label{tab:anpr}
\rowcolors{2}{rowgray}{white}
\begin{tabularx}{\textwidth}{lXcc}
\toprule
\textbf{Подход} & \textbf{Описание} & \textbf{mAP plate} & \textbf{Усл. освещения} \\
\midrule
Морфология + шаблоны & Пороговая обработка + шаблонный OCR & $<0.70$ & Плохо \\
HOG + SVM + Tesseract & HOG-дескриптор для локализации + Tesseract OCR & 0.75--0.80 & Удовлетворительно \\
YOLO + CRNN & YOLO-детекция + seq2seq OCR & 0.88--0.93 & Хорошо \\
\textbf{YOLOv8 + EasyOCR} & \textbf{Anchor-free детекция + CTC OCR} & \textbf{>0.90} & \textbf{Хорошо} \\
\bottomrule
\end{tabularx}
\end{table}

\noindent\textbf{Обоснование выбора:} YOLOv8n (fine-tuned) + EasyOCR обеспечивает
готовую поддержку кириллицы без обучения собственной OCR-модели.

\section{Обзор методов распознавания марки и типа кузова (VMMR)}

Распознавание марки/модели автомобиля (VMMR) относится к задачам тонкой
многоклассовой классификации~\cite{vmmr2024,siddiqui2016}. Задача:
\begin{equation}
  \hat{y} = \arg\max_{k}\; P(y = k \mid x),
\end{equation}
где $x$~--- изображение, $k$~--- класс марки или типа кузова.

\begin{table}[H]
\centering
\caption{Сравнение архитектур классификаторов для VMMR}
\label{tab:vmmr_compare}
\rowcolors{2}{rowgray}{white}
\begin{tabularx}{\textwidth}{lXccc}
\toprule
\textbf{Архитектура} & \textbf{Особенности} & \textbf{Params, M} & \textbf{Top-1 Acc} & \textbf{FPS} \\
\midrule
AlexNet        & Первая глубокая CNN для ImageNet     & 61   & 0.74--0.80 & высокий \\
VGG16          & Глубокая однородная архитектура      & 138  & 0.85--0.89 & низкий  \\
\textbf{ResNet18} & \textbf{Residual connections, лёгкая} & \textbf{11.7} & \textbf{0.88--0.92} & \textbf{высокий} \\
MobileNetV2    & Depthwise sep. conv., инв. бутылочное горлышко & 3.4 & 0.86--0.91 & очень высокий \\
EfficientNet-B0 & Составной масштаб (глубина/ширина/разрешение) & 5.3 & 0.90--0.94 & средний \\
\bottomrule
\end{tabularx}
\end{table}

\noindent\textbf{Обоснование выбора ResNet18:} оптимальное соотношение
точности и скорости; 11.7M параметров позволяют дообучить модель
на ограниченных данных без переобучения.

\section{Обзор методов определения цвета автомобиля}

\begin{table}[H]
\centering
\caption{Сравнение методов классификации цвета автомобиля}
\label{tab:color_methods}
\rowcolors{2}{rowgray}{white}
\begin{tabular}{llcc}
\toprule
\textbf{Метод} & \textbf{Принцип} & \textbf{Accuracy} & \textbf{Устойчивость к освещению} \\
\midrule
k-means (RGB)       & Кластеризация пикселей         & 0.70--0.81 & Низкая \\
k-means (HSV)       & Кластеризация в HSV-пространстве & 0.78--0.85 & Средняя \\
SVM + цветовые гистограммы & Гистограммы + SVM        & 0.83--0.88 & Средняя \\
\textbf{CNN (ResNet18)} & \textbf{Сквозное обучение}  & \textbf{>0.90} & \textbf{Высокая} \\
\bottomrule
\end{tabular}
\end{table}

Для повышения устойчивости в видеопотоке применяется агрегация предсказаний
по нескольким кадрам~\cite{kubatin2019}:
\begin{equation}
  \hat{y} = \arg\max_{k}\sum_{t=1}^{T} \mathbf{1}[y_t = k].
\end{equation}

\noindent\textbf{Обоснование выбора:} CNN (ResNet18) неявно учится
инвариантности к освещению и превосходит кластеризационные методы на 10--15\%.

\section{Цель и задачи работы}

\textbf{Цель}~--- разработка системы автоматического определения цвета,
типа кузова и номерного знака автомобиля по видеопотоку.

\textbf{Задачи:}
\begin{enumerate}
  \item Провести анализ существующих методов детекции ТС, ANPR, VMMR и определения цвета.
  \item Выбрать и обосновать архитектуры моделей для каждой подзадачи.
  \item Подготовить обучающие датасеты.
  \item Обучить и оценить модели детекции номеров, классификаторы цвета и типа кузова.
  \item Интегрировать модели в единый конвейер обработки видео.
  \item Провести сравнительный эксперимент ResNet18 vs MobileNetV2.
\end{enumerate}

%% =============================================================
%%  ГЛАВА 2 — МЕТОДЫ И ДАТАСЕТЫ
%% =============================================================
\chapter{Описание выбранных методов и датасетов}

\section{Архитектура системы}

Разработанная система представляет собой последовательный конвейер,
обрабатывающий каждый кадр видеопотока. Общая схема приведена
на рисунке~\ref{fig:pipeline}.

\begin{figure}[H]
\centering
\begin{tikzpicture}[node distance=0.55cm and 0.6cm]

  %% ── левая колонка (основной поток) ──────────────────────
  \node[blk] (frame) {Кадр\\видео};
  \node[blk, below=of frame] (vdet)  {YOLOv8n\\Детекция ТС};
  \node[blk, below=of vdet]  (crop)  {Кроп\\автомобиля};
  \node[blk, below=of crop]  (pdet)  {YOLOv8n\\Детекция\\номера};
  \node[blk, below=of pdet]  (ocr)   {EasyOCR\\Чтение\\номера};
  \node[grn, below=of ocr]   (out)   {Результат\\+ визуализация};

  \draw[arr] (frame) -- (vdet);
  \draw[arr] (vdet)  -- node[right, font=\scriptsize]{bbox} (crop);
  \draw[arr] (crop)  -- (pdet);
  \draw[arr] (pdet)  -- node[right, font=\scriptsize]{кроп номера} (ocr);
  \draw[arr] (ocr)   -- (out);

  %% ── правая ветка (классификаторы) ───────────────────────
  \node[org, right=2.5cm of crop] (color) {ResNet18\\Цвет\\(9 классов)};
  \node[org, below=0.55cm of color] (body)  {ResNet18\\Кузов\\(6 классов)};

  \draw[arr] (crop.east) -- ++(0.6,0) |- (color.west);
  \draw[arr] (crop.east) -- ++(0.6,0) |- (body.west);
  \draw[arr] (color.south) |- (out.east);
  \draw[arr] (body.east)  -- ++(0.3,0) |- (out.east);

  %% ── подпись классов ─────────────────────────────────────
  \node[gry, right=2.5cm of color, text width=2.8cm, font=\scriptsize]
        (clbl) {black, white, silver,\\red, blue, yellow,\\green, brown, orange};
  \node[gry, right=2.5cm of body,  text width=2.8cm, font=\scriptsize]
        (blbl) {sedan, suv,\\hatchback, van,\\truck, minibus};
  \draw[darr] (color.east) -- (clbl.west);
  \draw[darr] (body.east)  -- (blbl.west);

\end{tikzpicture}
\caption{Схема конвейера обработки видеокадра}
\label{fig:pipeline}
\end{figure}

\section{Детекция транспортных средств: YOLOv8}

\subsection{Внутренняя архитектура YOLOv8}

YOLOv8 состоит из трёх функциональных блоков, показанных
на рисунке~\ref{fig:yolo_arch}.

\begin{figure}[H]
\centering
\begin{tikzpicture}[node distance=0.5cm and 0.8cm, every node/.style={font=\small}]

  \node[blk, text width=2.6cm] (inp) {Входное изображение\\$640\times640$};

  \node[blk, text width=3cm, right=0.9cm of inp] (bb)
        {\textbf{Backbone}\\CSPDarkNet-v8\\(C2f блоки)};

  \node[blk, text width=3cm, right=0.9cm of bb] (neck)
        {\textbf{Neck}\\PAN-FPN\\мультимасштаб};

  \node[grn, text width=3cm, right=0.9cm of neck] (head)
        {\textbf{Head}\\Anchor-free\\Bbox + Cls};

  \node[org, text width=3cm, right=0.9cm of head] (out)
        {Детекции\\$(x,y,w,h,p,c)$};

  \draw[arr] (inp)  -- (bb);
  \draw[arr] (bb)   -- node[above,font=\scriptsize]{P3/P4/P5} (neck);
  \draw[arr] (neck) -- (head);
  \draw[arr] (head) -- (out);

  %% подписи снизу
  \node[below=0.5cm of bb,   font=\scriptsize, text=graynode]
       {Feature extraction};
  \node[below=0.5cm of neck, font=\scriptsize, text=graynode]
       {Feature fusion};
  \node[below=0.5cm of head, font=\scriptsize, text=graynode]
       {Prediction};

\end{tikzpicture}
\caption{Упрощённая схема архитектуры YOLOv8}
\label{fig:yolo_arch}
\end{figure}

Функция потерь YOLOv8:
\begin{equation}
  \mathcal{L} = \lambda_{\mathrm{box}}\,\mathcal{L}_{\mathrm{box}}
              + \lambda_{\mathrm{cls}}\,\mathcal{L}_{\mathrm{cls}}
              + \lambda_{\mathrm{dfl}}\,\mathcal{L}_{\mathrm{dfl}},
\end{equation}
где $\mathcal{L}_{\mathrm{box}}$~--- CIoU loss,
$\mathcal{L}_{\mathrm{cls}}$~--- BCE,
$\mathcal{L}_{\mathrm{dfl}}$~--- Distribution Focal Loss.

\subsection{Параметры инференса}

\begin{table}[H]
\centering
\caption{Параметры запуска YOLOv8n для детекции ТС}
\label{tab:yolo_params}
\begin{tabular}{ll}
\toprule
\textbf{Параметр} & \textbf{Значение} \\
\midrule
Модель                  & \texttt{yolov8n.pt} (предобученная, COCO) \\
Порог уверенности       & $p_{\min} = 0.40$ \\
Порог NMS IoU           & $0.45$ \\
Используемые классы     & 2 (car), 3 (motorcycle), 5 (bus), 7 (truck) \\
Размер входа            & $640 \times 640$ \\
Аппаратная платформа    & GPU NVIDIA Tesla T4 / CPU \\
\bottomrule
\end{tabular}
\end{table}

\section{Детекция номерных знаков: дообученный YOLOv8n}

\subsection{Схема transfer learning}

Процесс дообучения показан на рисунке~\ref{fig:transfer}.

\begin{figure}[H]
\centering
\begin{tikzpicture}[node distance=0.6cm and 0.7cm]

  \node[blk, text width=3.2cm] (pre)
        {YOLOv8n\\предобученная\\(COCO, 80 классов)};

  \node[blk, text width=3.2cm, right=1.2cm of pre] (frz)
        {Заморозка\\Backbone\\(первые эпохи)};

  \node[org, text width=3.2cm, right=1.2cm of frz] (head)
        {Переобучение\\Head\\(1 класс: plate)};

  \node[grn, text width=3.2cm, right=1.2cm of head] (res)
        {Финальная модель\\yolov8n\_plates.pt};

  \draw[arr] (pre)  -- (frz);
  \draw[arr] (frz)  -- (head);
  \draw[arr] (head) -- (res);

  \node[below=0.7cm of frz, font=\scriptsize, text width=3.8cm, align=center]
       {Roboflow License Plates\\$\approx$6000 изображений};
  \draw[darr] ([yshift=-0.5cm]frz.south) -- (head.south |- {[yshift=-0.5cm]frz.south});

\end{tikzpicture}
\caption{Схема transfer learning для детектора номерных знаков}
\label{fig:transfer}
\end{figure}

\subsection{Параметры обучения}

\begin{table}[H]
\centering
\caption{Гиперпараметры дообучения детектора номеров}
\label{tab:plate_hp}
\begin{tabular}{ll}
\toprule
\textbf{Параметр} & \textbf{Значение} \\
\midrule
Базовая модель       & \texttt{yolov8n.pt} \\
Эпохи                & 50 \\
Размер входа         & $640 \times 640$ \\
Размер батча         & 16 \\
Оптимизатор          & SGD \\
Начальная LR         & 0.01 \\
Momentum             & 0.937 \\
Weight decay         & $5 \times 10^{-4}$ \\
Число классов        & 1 (\textit{license\_plate}) \\
Метрика сохранения   & mAP@0.50 на валидации \\
\bottomrule
\end{tabular}
\end{table}

\section{Оптическое распознавание номеров: EasyOCR}

EasyOCR использует двухэтапный конвейер, представленный
на рисунке~\ref{fig:ocr_pipeline}.

\begin{figure}[H]
\centering
\begin{tikzpicture}[node distance=0.5cm and 0.7cm]
  \node[blk, text width=3cm]  (plate) {Кроп номерного\\знака (BGR)};
  \node[blk, text width=3cm, right=0.9cm of plate]  (gray)  {Оттенки\\серого\\(Otsu)};
  \node[blk, text width=3.2cm, right=0.9cm of gray]  (craft) {\textbf{CRAFT}\\Карты символов\\+ аффинности};
  \node[blk, text width=3.2cm, right=0.9cm of craft] (crnn)  {\textbf{CRNN+CTC}\\Декодирование\\строки};
  \node[grn, text width=3cm, right=0.9cm of crnn]   (regex) {Regex-фильтр\\RU формата};

  \draw[arr] (plate) -- (gray);
  \draw[arr] (gray)  -- (craft);
  \draw[arr] (craft) -- node[above,font=\scriptsize]{регионы} (crnn);
  \draw[arr] (crnn)  -- (regex);
\end{tikzpicture}
\caption{Конвейер EasyOCR для распознавания номерного знака}
\label{fig:ocr_pipeline}
\end{figure}

Качество OCR оценивается через Character Error Rate:
\begin{equation}
  \mathrm{CER} = \frac{S + D + I}{N},
\end{equation}
где $S, D, I$~--- замены, удаления, вставки (алгоритм Левенштейна),
$N$~--- длина эталонной строки.

\section{Классификация цвета и типа кузова: ResNet18}

\subsection{Архитектура ResNet18}

Ключевая идея ResNet~--- shortcut-соединения:
\begin{equation}
  \mathbf{y} = \mathcal{F}(\mathbf{x},\{W_i\}) + \mathbf{x},
\end{equation}
что устраняет затухание градиентов и позволяет строить глубокие сети.
Схема residual-блока и размещение слоёв в ResNet18 приведены
на рисунке~\ref{fig:resnet_arch}.

\begin{figure}[H]
\centering
\begin{tikzpicture}[node distance=0.45cm and 0.8cm]

  %% ── Residual-блок ────────────────────────────────────────
  \begin{scope}[local bounding box=resblk]
    \node[blk, text width=2.4cm] (in)   {Вход $\mathbf{x}$};
    \node[blk, text width=2.4cm, below=of in]   (c1) {Conv $3\times3$\\BN + ReLU};
    \node[blk, text width=2.4cm, below=of c1]   (c2) {Conv $3\times3$\\BN};
    \node[grn, text width=2.4cm, below=of c2]   (add){$+$ (shortcut)\\ReLU};
    \node[blk, text width=2.4cm, below=of add]  (out){Выход $\mathbf{y}$};

    \draw[arr] (in)  -- (c1);
    \draw[arr] (c1)  -- (c2);
    \draw[arr] (c2)  -- (add);
    \draw[arr] (add) -- (out);

    %% shortcut
    \draw[arr, bluenode, rounded corners=6pt]
        (in.east) -- ++(1.0,0) -- ++(0,-4.0) -- (add.east);
    \node[right=1.15cm of c1, font=\scriptsize, text=bluenode]
         {shortcut};
  \end{scope}

  %% ── Состав ResNet18 ──────────────────────────────────────
  \begin{scope}[xshift=6.5cm, node distance=0.35cm]
    \node[blk, text width=3.4cm] (r0)  {Conv $7\times7$, s=2\\BN + ReLU};
    \node[blk, text width=3.4cm, below=of r0]  (mp)  {MaxPool $3\times3$, s=2};
    \node[blk, text width=3.4cm, below=of mp]  (r1)  {Layer1: $2\times$ Res-Block\\64 каналов};
    \node[blk, text width=3.4cm, below=of r1]  (r2)  {Layer2: $2\times$ Res-Block\\128 каналов};
    \node[blk, text width=3.4cm, below=of r2]  (r3)  {Layer3: $2\times$ Res-Block\\256 каналов};
    \node[blk, text width=3.4cm, below=of r3]  (r4)  {Layer4: $2\times$ Res-Block\\512 каналов};
    \node[blk, text width=3.4cm, below=of r4]  (gap) {AvgPool + Flatten};
    \node[grn, text width=3.4cm, below=of gap] (fc)  {FC: 512 → $K$ классов};

    \foreach \a/\b in {r0/mp, mp/r1, r1/r2, r2/r3, r3/r4, r4/gap, gap/fc}
      \draw[arr] (\a) -- (\b);
  \end{scope}

  %% подпись блока
  \node[above=0.1cm of in, font=\small\bfseries] {Residual-блок};
  \node[above=0.1cm of r0, font=\small\bfseries] {Состав ResNet18};

\end{tikzpicture}
\caption{Residual-блок и структура ResNet18 (18 слоёв, $K$~--- число классов)}
\label{fig:resnet_arch}
\end{figure}

\subsection{Двухфазное дообучение}

Стратегия двухфазного обучения показана на рисунке~\ref{fig:twophase}.

\begin{figure}[H]
\centering
\begin{tikzpicture}[node distance=0.6cm and 1.0cm]
  \node[blk, text width=4cm] (p1)
        {\textbf{Фаза 1}\\Backbone заморожен\\Обучается только FC\\$\eta_1 = 10^{-3}$};
  \node[org, text width=3cm, right=1.2cm of p1] (sw)
        {Эпоха\\UNFREEZE\_AT\\Разморозка};
  \node[blk, text width=4cm, right=1.2cm of sw] (p2)
        {\textbf{Фаза 2}\\Весь backbone обучается\\$\eta_2 = \eta_1/10$};
  \node[grn, text width=3cm, right=1.2cm of p2] (save)
        {Best model\\(max val acc)};

  \draw[arr] (p1) -- node[above,font=\scriptsize]{эп. 1–N} (sw);
  \draw[arr] (sw) -- node[above,font=\scriptsize]{эп. N+1–end} (p2);
  \draw[arr] (p2) -- (save);
\end{tikzpicture}
\caption{Стратегия двухфазного дообучения ResNet18}
\label{fig:twophase}
\end{figure}

Функция потерь~--- кросс-энтропия:
\begin{equation}
  \mathcal{L}_{\mathrm{CE}} = -\sum_{k=1}^{K} y_k \log \hat{y}_k,
\end{equation}
где $y_k$~--- one-hot метка, $\hat{y}_k$~--- предсказанная вероятность.

\subsection{Гиперпараметры обучения}

\begin{table}[H]
\centering
\caption{Гиперпараметры обучения классификаторов ResNet18}
\label{tab:resnet_hp}
\rowcolors{2}{rowgray}{white}
\begin{tabular}{lcc}
\toprule
\textbf{Параметр} & \textbf{Классификатор цвета} & \textbf{Классификатор кузова} \\
\midrule
Число классов     & 9              & 6 \\
Эпохи             & 20             & 25 \\
Разморозка на эпохе & 5            & 6 \\
Батч              & 32             & 32 \\
LR (Фаза 1)       & $10^{-3}$      & $10^{-3}$ \\
LR (Фаза 2)       & $10^{-4}$      & $10^{-4}$ \\
Оптимизатор       & Adam           & Adam \\
$\beta_1, \beta_2$ & 0.9, 0.999   & 0.9, 0.999 \\
StepLR шаг        & 5 эпох         & 6 эпох \\
StepLR $\gamma$   & 0.3            & 0.3 \\
Размер входа      & $224\times224$ & $224\times224$ \\
\bottomrule
\end{tabular}
\end{table}

\subsection{Аугментация данных}

\begin{table}[H]
\centering
\caption{Аугментации для обучения классификаторов}
\label{tab:augmentation}
\begin{tabular}{lll}
\toprule
\textbf{Преобразование} & \textbf{Параметры} & \textbf{Назначение} \\
\midrule
\multicolumn{3}{l}{\textit{Обучающая выборка}} \\
RandomResizedCrop   & $224$, scale $[0.65, 1.0]$    & Масштаб и позиция \\
RandomHorizontalFlip & $p=0.5$                      & Зеркальное отражение \\
ColorJitter         & brightness=0.3, contrast=0.3  & Освещение \\
RandomRotation      & $\pm 10°$                     & Поворот \\
Normalize           & $\mu=[.485,.456,.406]$         & ImageNet stats \\
\midrule
\multicolumn{3}{l}{\textit{Валидационная выборка}} \\
Resize              & $224\times224$                 & Размер входа \\
Normalize           & $\mu=[.485,.456,.406]$         & ImageNet stats \\
\bottomrule
\end{tabular}
\end{table}

\section{Сравниваемая архитектура: MobileNetV2}

\begin{figure}[H]
\centering
\begin{tikzpicture}[node distance=0.5cm and 0.9cm]

  %% Inverted Residual Block
  \begin{scope}
    \node[blk, text width=3cm] (mi)  {Вход $h \times w \times c$};
    \node[blk, text width=3cm, below=of mi] (exp) {Conv $1\times1$\\(расширение $\times t$)};
    \node[blk, text width=3cm, below=of exp](dw)  {DWConv $3\times3$\\(depthwise)};
    \node[blk, text width=3cm, below=of dw] (proj){Conv $1\times1$\\(проекция)};
    \node[grn, text width=3cm, below=of proj](mo) {Выход $h \times w \times c'$};

    \draw[arr] (mi)  -- (exp);
    \draw[arr] (exp) -- (dw);
    \draw[arr] (dw)  -- (proj);
    \draw[arr] (proj)-- (mo);

    \draw[arr, orangenode, rounded corners=6pt]
        (mi.east) -- ++(0.9,0) -- ++(0,-3.7) -- (mo.east);
    \node[right=1.05cm of dw, font=\scriptsize, text=orangenode]{skip (если $c=c'$)};
  \end{scope}

  %% Сравнение параметров
  \begin{scope}[xshift=6.5cm]
    \node[font=\small\bfseries] at (1.7, 0) {ResNet18 vs MobileNetV2};
    \node[gry, text width=4.8cm, below=0.3cm of {(1.7,0)}] (tbl)
    {
    \begin{tabular}{lcc}
    \small Характеристика & \small R18 & \small MV2 \\
    \hline
    \small Params, M       & \small 11.7 & \small 3.4 \\
    \small Blocks          & \small Residual & \small Inv.Res. \\
    \small Conv type       & \small Standard & \small Depthwise \\
    \small FLOPS (224)     & \small 1.8G & \small 0.3G \\
    \end{tabular}
    };
  \end{scope}

  \node[above=0.1cm of mi, font=\small\bfseries] {Inverted Residual Block};
\end{tikzpicture}
\caption{Инвертированный residual-блок MobileNetV2 и сравнение с ResNet18}
\label{fig:mobilenet}
\end{figure}

\section{Датасеты}

\begin{table}[H]
\centering
\caption{Сводная информация об используемых датасетах}
\label{tab:datasets_summary}
\rowcolors{2}{rowgray}{white}
\begin{tabularx}{\textwidth}{lXcccc}
\toprule
\textbf{Датасет} & \textbf{Задача} & \textbf{Источник} & \textbf{Изобр.} & \textbf{Классов} & \textbf{Формат} \\
\midrule
COCO 2017          & Детекция ТС        & COCO/Ultralytics & 118k & 80 (4 для ТС) & YOLO \\
Roboflow LP        & Детекция номеров   & Roboflow         & $\approx$6k  & 1 & YOLO \\
VCoR               & Цвет авто          & Kaggle           & $\approx$15k & 9 & ImageFolder \\
CompCars body-type & Тип кузова         & CUHK             & $\approx$30k & 6 & ImageFolder \\
\bottomrule
\end{tabularx}
\end{table}

\subsection{VCoR~--- классификация цвета}

\begin{table}[H]
\centering
\caption{Классы датасета VCoR и соответствующие цвета}
\label{tab:vcor}
\rowcolors{2}{rowgray}{white}
\begin{tabular}{clll}
\toprule
\textbf{\#} & \textbf{Класс} & \textbf{Приблизительный RGB} & \textbf{Hex} \\
\midrule
1 & black  & (25, 25, 25)   & \#191919 \\
2 & white  & (230, 230, 230)& \#E6E6E6 \\
3 & silver & (175, 180, 185)& \#AFB4B9 \\
4 & red    & (195, 30, 30)  & \#C31E1E \\
5 & blue   & (30, 60, 200)  & \#1E3CC8 \\
6 & yellow & (220, 200, 20) & \#DCC814 \\
7 & green  & (30, 145, 30)  & \#1E911E \\
8 & brown  & (100, 60, 30)  & \#643C1E \\
9 & orange & (215, 105, 20) & \#D7691E \\
\bottomrule
\end{tabular}
\end{table}

\subsection{CompCars~--- классификация типа кузова}

\begin{table}[H]
\centering
\caption{Маппинг оригинальных классов CompCars на целевые}
\label{tab:compcars}
\rowcolors{2}{rowgray}{white}
\begin{tabular}{lll}
\toprule
\textbf{Оригинальный класс CompCars} & \textbf{Целевой класс} & \textbf{Описание} \\
\midrule
sedan, fastback, convertible & sedan    & Легковой, 4 двери \\
hatchback, estate            & hatchback& Хетчбэк, универсал \\
SUV, crossover               & suv      & Внедорожник \\
MPV                          & van      & Минивэн \\
pickup, truck                & truck    & Грузовой / пикап \\
microbus, minibus            & minibus  & Микроавтобус \\
\bottomrule
\end{tabular}
\end{table}

%% =============================================================
%%  ГЛАВА 3 — ЭКСПЕРИМЕНТЫ
%% =============================================================
\chapter{Описание экспериментов}

\section{Условия проведения экспериментов}

\begin{table}[H]
\centering
\caption{Программно-аппаратная конфигурация экспериментов}
\label{tab:env}
\begin{tabular}{ll}
\toprule
\textbf{Компонент} & \textbf{Значение} \\
\midrule
Платформа         & Google Colab \\
GPU               & NVIDIA Tesla T4, 16 GB VRAM \\
CPU               & Intel Xeon (2 vCPU) \\
RAM               & 12 GB \\
OS                & Ubuntu 20.04 \\
Python            & 3.10 \\
PyTorch           & 2.x + CUDA 11.8 \\
Ultralytics       & 8.x \\
EasyOCR           & 1.7 \\
torchvision       & 0.17 \\
\bottomrule
\end{tabular}
\end{table}

Ноутбуки для воспроизведения:
\begin{itemize}[nosep]
  \item \texttt{01\_detection.ipynb}~--- детекция ТС, fine-tune детектора номеров, OCR;
  \item \texttt{02\_color\_classifier.ipynb}~--- классификатор цвета (VCoR);
  \item \texttt{03\_body\_classifier.ipynb}~--- классификатор кузова (CompCars);
  \item \texttt{04\_pipeline\_final.ipynb}~--- финальный конвейер, статистика.
\end{itemize}

\section{Эксперимент 1: дообучение детектора номерных знаков}

\subsection{Протокол}

YOLOv8n дообучался на датасете Roboflow License Plates в течение 50 эпох
(параметры~--- таблица~\ref{tab:plate_hp}). Лучшая модель сохранялась
по максимуму mAP@0.50 на валидации.

\subsection{Результаты}

\begin{table}[H]
\centering
\caption{Метрики детектора номерных знаков (val)}
\label{tab:plate_det}
\rowcolors{2}{rowgray}{white}
\begin{tabular}{lcc}
\toprule
\textbf{Метрика} & \textbf{YOLOv8n (pretrained)} & \textbf{YOLOv8n (fine-tuned)} \\
\midrule
mAP@0.50       & --- & \textit{(заполнить)} \\
mAP@0.50:0.95  & --- & \textit{(заполнить)} \\
Precision      & --- & \textit{(заполнить)} \\
Recall         & --- & \textit{(заполнить)} \\
\bottomrule
\end{tabular}
\end{table}

% Кривые обучения:
% \begin{figure}[H]
%   \centering
%   \includegraphics[width=0.85\textwidth]{figs/plate_training_curves.png}
%   \caption{Кривые loss и mAP при дообучении детектора номеров}
%   \label{fig:plate_curves}
% \end{figure}

\section{Эксперимент 2: классификатор цвета}

\subsection{Протокол}

ResNet18 (ImageNet) дообучался на VCoR, 20 эпох, разморозка на эпохе~5.
Батч: 32. Аугментации~--- таблица~\ref{tab:augmentation}.

\subsection{Результаты}

\begin{table}[H]
\centering
\caption{Метрики классификатора цвета (val)}
\label{tab:color_metrics}
\rowcolors{2}{rowgray}{white}
\begin{tabular}{lcc}
\toprule
\textbf{Метрика} & \textbf{Baseline (fc only)} & \textbf{Fine-tuned (2 phase)} \\
\midrule
Accuracy top-1 & \textit{(заполнить)} & \textit{(заполнить)} \\
Accuracy top-3 & \textit{(заполнить)} & \textit{(заполнить)} \\
Macro F1       & \textit{(заполнить)} & \textit{(заполнить)} \\
\bottomrule
\end{tabular}
\end{table}

\begin{table}[H]
\centering
\caption{F1-score по классам цвета (val, fine-tuned модель)}
\label{tab:color_per_class}
\rowcolors{2}{rowgray}{white}
\begin{tabular}{lccc}
\toprule
\textbf{Класс} & \textbf{Precision} & \textbf{Recall} & \textbf{F1} \\
\midrule
black   & \textit{---} & \textit{---} & \textit{---} \\
white   & \textit{---} & \textit{---} & \textit{---} \\
silver  & \textit{---} & \textit{---} & \textit{---} \\
red     & \textit{---} & \textit{---} & \textit{---} \\
blue    & \textit{---} & \textit{---} & \textit{---} \\
yellow  & \textit{---} & \textit{---} & \textit{---} \\
green   & \textit{---} & \textit{---} & \textit{---} \\
brown   & \textit{---} & \textit{---} & \textit{---} \\
orange  & \textit{---} & \textit{---} & \textit{---} \\
\midrule
\textbf{Macro avg} & \textit{---} & \textit{---} & \textit{---} \\
\bottomrule
\end{tabular}
\end{table}

% Confusion matrix и F1-bars:
% \begin{figure}[H]
%   \centering
%   \begin{subfigure}[b]{0.48\textwidth}
%     \includegraphics[width=\textwidth]{figs/color_confusion_matrix.png}
%     \caption{Confusion matrix}
%   \end{subfigure}
%   \hfill
%   \begin{subfigure}[b]{0.48\textwidth}
%     \includegraphics[width=\textwidth]{figs/color_per_class_metrics.png}
%     \caption{F1 по классам}
%   \end{subfigure}
%   \caption{Оценка классификатора цвета}
%   \label{fig:color_eval}
% \end{figure}

\section{Эксперимент 3: классификатор типа кузова}

\subsection{Протокол}

ResNet18 дообучался на CompCars body type, 25 эпох, разморозка на эпохе~6.
Параметры аналогичны эксперименту~2.

\subsection{Результаты}

\begin{table}[H]
\centering
\caption{Метрики классификатора типа кузова (val)}
\label{tab:body_metrics}
\rowcolors{2}{rowgray}{white}
\begin{tabular}{lcc}
\toprule
\textbf{Метрика} & \textbf{Baseline (fc only)} & \textbf{Fine-tuned (2 phase)} \\
\midrule
Accuracy top-1 & \textit{(заполнить)} & \textit{(заполнить)} \\
Accuracy top-3 & \textit{(заполнить)} & \textit{(заполнить)} \\
Macro F1       & \textit{(заполнить)} & \textit{(заполнить)} \\
\bottomrule
\end{tabular}
\end{table}

\begin{table}[H]
\centering
\caption{F1-score по классам типа кузова (val, fine-tuned модель)}
\label{tab:body_per_class}
\rowcolors{2}{rowgray}{white}
\begin{tabular}{lccc}
\toprule
\textbf{Класс} & \textbf{Precision} & \textbf{Recall} & \textbf{F1} \\
\midrule
sedan     & \textit{---} & \textit{---} & \textit{---} \\
suv       & \textit{---} & \textit{---} & \textit{---} \\
hatchback & \textit{---} & \textit{---} & \textit{---} \\
van       & \textit{---} & \textit{---} & \textit{---} \\
truck     & \textit{---} & \textit{---} & \textit{---} \\
minibus   & \textit{---} & \textit{---} & \textit{---} \\
\midrule
\textbf{Macro avg} & \textit{---} & \textit{---} & \textit{---} \\
\bottomrule
\end{tabular}
\end{table}

\section{Эксперимент 4: сравнение архитектур ResNet18 и MobileNetV2}

\subsection{Постановка}

Оба классификатора обучались на одинаковом разбиении CompCars при одинаковых
гиперпараметрах (только голова $\mathrm{fc}$, без разморозки backbone).
Измерялись val Accuracy и FPS инференса ($224 \times 224$, GPU~T4).

\subsection{Результаты}

\begin{table}[H]
\centering
\caption{Сравнение архитектур на задаче классификации типа кузова}
\label{tab:arch_comparison}
\rowcolors{2}{rowgray}{white}
\begin{tabular}{lccccc}
\toprule
\textbf{Модель} & \textbf{Params, M} & \textbf{FLOPS, G} & \textbf{Val Acc} & \textbf{FPS (T4)} & \textbf{Размер .pth, MB} \\
\midrule
ResNet18    & 11.7 & 1.8 & \textit{(заполнить)} & \textit{(заполнить)} & $\approx 45$ \\
MobileNetV2 &  3.4 & 0.3 & \textit{(заполнить)} & \textit{(заполнить)} & $\approx 14$ \\
\bottomrule
\end{tabular}
\end{table}

\section{Эксперимент 5: финальный конвейер}

\subsection{Тест на одиночном изображении}

Для качественной проверки конвейер запускается на тестовом изображении
(Toyota Camry 2011, Wikimedia Commons). Ожидаемый результат: детекция
автомобиля, предсказание цвета \textit{silver}, типа кузова \textit{sedan},
распознавание номера.

% \begin{figure}[H]
%   \centering
%   \includegraphics[width=0.75\textwidth]{figs/test_single_frame.jpg}
%   \caption{Результат обработки тестового изображения финальным конвейером}
%   \label{fig:test_frame}
% \end{figure}

\subsection{Метрики производительности}

\begin{table}[H]
\centering
\caption{Задержка и производительность модулей конвейера (GPU~T4)}
\label{tab:pipeline_fps}
\rowcolors{2}{rowgray}{white}
\begin{tabular}{lcc}
\toprule
\textbf{Модуль} & \textbf{Задержка, мс} & \textbf{Доля в общей задержке, \%} \\
\midrule
YOLOv8n детекция ТС       & \textit{(заполнить)} & \textit{---} \\
YOLOv8n детекция номера   & \textit{(заполнить)} & \textit{---} \\
ResNet18 цвет             & \textit{(заполнить)} & \textit{---} \\
ResNet18 кузов            & \textit{(заполнить)} & \textit{---} \\
EasyOCR                   & \textit{(заполнить)} & \textit{---} \\
\midrule
\textbf{Итого (pipeline)} & \textit{(заполнить)} & 100\% \\
\textbf{FPS}              & \multicolumn{2}{c}{\textit{(заполнить)}} \\
\bottomrule
\end{tabular}
\end{table}

\subsection{Итоговая таблица результатов}

\begin{table}[H]
\centering
\caption{Сводная таблица метрик по всем модулям}
\label{tab:final_summary}
\rowcolors{2}{rowgray}{white}
\begin{tabular}{llcl}
\toprule
\textbf{Модуль} & \textbf{Компонент} & \textbf{Целевая метрика} & \textbf{Результат} \\
\midrule
Детекция ТС         & YOLOv8n (COCO)           & mAP@0.50        & $>0.80$ (pretrained) \\
Детекция номеров    & YOLOv8n (fine-tuned)     & mAP@0.50        & \textit{(заполнить)} \\
Классификация цвета & ResNet18 (VCoR)          & Accuracy top-1  & \textit{(заполнить)} \\
Классификация кузова& ResNet18 (CompCars)      & Accuracy top-1  & \textit{(заполнить)} \\
OCR номеров         & EasyOCR ru+en            & CER             & \textit{(заполнить)} \\
\textbf{Pipeline}   & \textbf{Полный конвейер} & \textbf{FPS}    & \textit{(заполнить)} \\
\bottomrule
\end{tabular}
\end{table}

%% =============================================================
%%  ГЛАВА 4 — ВЫВОДЫ
%% =============================================================
\chapter{Выводы}

В рамках выпускной квалификационной работы разработана система автоматического
определения характеристик транспортных средств. Сформулированы следующие выводы.

\begin{enumerate}
  \item \textbf{Архитектурный выбор.}
        Комбинация YOLOv8n~+ ResNet18~+ EasyOCR обеспечивает приемлемый баланс
        точности и скорости (см. таблицу~\ref{tab:final_summary}).

  \item \textbf{Детектор номерных знаков.}
        Fine-tuning YOLOv8n на датасете Roboflow LP позволяет достичь
        mAP@0.50\,$>$\,\textit{(XX)}, что подтверждает эффективность
        transfer learning при малом объёме целевых данных.

  \item \textbf{Классификация цвета.}
        ResNet18, дообученный на VCoR за 20 эпох с двухфазной стратегией,
        достигает Accuracy top-1\,$=$\,\textit{(XX)}, превосходя
        кластеризационные методы на 10--15\,\%
        (таблица~\ref{tab:color_methods}).

  \item \textbf{Классификация типа кузова.}
        Модель достигает Accuracy top-1\,$=$\,\textit{(XX)}.
        Сравнение ResNet18 и MobileNetV2 (таблица~\ref{tab:arch_comparison})
        показывает, что \textit{(вставить вывод по результатам)}.

  \item \textbf{Производительность конвейера.}
        Полный конвейер обрабатывает видеопоток со скоростью \textit{(XX)} FPS
        на GPU~T4. Основная задержка~--- EasyOCR (\textit{XX}\,\% от общего времени).

  \item \textbf{Ограничения системы.}
        EasyOCR является узким местом конвейера. Для встроенных систем
        целесообразно рассмотреть специализированный лёгкий OCR-модуль.
\end{enumerate}

\textbf{Направления дальнейшего развития:}
\begin{itemize}[nosep]
  \item замена EasyOCR на CRNN-детектор, оптимизированный для российских номеров;
  \item расширение классов цвета и типа кузова;
  \item адаптация к ночным условиям съёмки (low-light enhancement);
  \item квантизация и pruning для развёртывания на NVIDIA Jetson.
\end{itemize}

%% =====================  ЛИТЕРАТУРА  =========================
\begin{thebibliography}{99}

\bibitem{laroca2018}
Laroca R., Severo E., Zanlorensi L.\,A. и др.
\textit{A Robust Real-Time Automatic License Plate Recognition Based on the YOLO Detector}~//
arXiv:1802.09567, 2018.

\bibitem{laroca2021}
Laroca R.
\textit{Automatic License Plate Recognition in Real-World Scenarios}: PhD Thesis~//
Universidade Federal do Paraná, 2021.

\bibitem{abramov2019}
Абрамов А.\,А.
\textit{Алгоритм распознавания номерных знаков в условиях плохого качества видео}~//
КиберЛенинка, 2019.

\bibitem{survey2021}
\textit{A Survey of Automatic License Plate Recognition Systems}~//
Sensors, 21(9):3028, 2021.

\bibitem{vmmr2024}
\textit{Two decades of vehicle make and model recognition}~//
Journal of King Saud University~--- CIS, 2024.

\bibitem{siddiqui2016}
Siddiqui A., Mammeri A., Boukerche A.
\textit{Real-Time Vehicle Make and Model Recognition Based on SURF Features}~//
IEEE ITS, 2016.

\bibitem{manzoor2022}
Manzoor M.\,A., Morgan Y.
\textit{Vehicle Make and Model Recognition Using Deep Learning}~//
Machine Learning and Knowledge Extraction, 1(2):36, 2022.

\bibitem{kubatin2019}
Кубатин И.\,А.
\textit{Алгоритм определения цвета автомобиля методом кластеризации k-средних}~//
КиберЛенинка, 2019.

\bibitem{wu2013}
Wu Y., Lim J., Yang M.-H.
\textit{Online Object Tracking: A Benchmark}~//
CVPR, 2013.

\bibitem{su2015}
Su H., Wu X., Chen S., Ji Q.
\textit{Vehicle Color Recognition Using Convolutional Neural Network}~//
IEEE ICIP, 2015.

\end{thebibliography}

\end{document}
